\subsection{Pseudometrics and metric quotients}\label{subsec:preliminaries-quotients}
\paraid{p-dcc3379498b8}
We prove estimates that also apply when distinct points are identified.

\paraid{p-f455c68e0d8f}
The diameter is continuous on
$\GHSpace$,
by \cite[arXiv v1, Example~6.29]{Tuzhilin2020},
which states that for all nonempty compact metric spaces
$\MetricSpace{\BaseCarrier}{\MetricSymbol}$
and
$\MetricSpace{\ComparisonCarrier}{\ComparisonMetric}$
we have
\[
 \abs{\diam\MetricSpace{\BaseCarrier }{\MetricSymbol }-\diam\MetricSpace{\ComparisonCarrier }{\ComparisonMetric }}\le 2\GHDistance(\MetricSpace{\BaseCarrier }{\MetricSymbol },\MetricSpace{\ComparisonCarrier }{\ComparisonMetric }).
\]
A continuous pseudometric
$\Pseudometric $
on a compact space
$\BaseCarrier $
defines an equivalence relation as follows.
For
$\BasePoint,\ComparisonPoint\in\BaseCarrier$,
write
$\BasePoint\sim_\Pseudometric\ComparisonPoint$
if
$\Pseudometric(\BasePoint,\ComparisonPoint)=0$.
Its metric quotient is
\[
 \MetricQuotient{\BaseCarrier }{\Pseudometric }
 =\MetricSpace{\BaseCarrier /{\sim_\Pseudometric }}{\overline{\Pseudometric }},
\]
where for representatives
$\BasePoint,\ComparisonPoint\in\BaseCarrier$
the quotient metric satisfies
\[
 \overline{\Pseudometric }([\BasePoint ],[\ComparisonPoint ])=\Pseudometric (\BasePoint ,\ComparisonPoint ).
\]
This quotient is compact.
When
$\Pseudometric $
is a metric,
it is isometric to
$\MetricSpace{\BaseCarrier }{\Pseudometric }$.
For the comparison of a pseudometric quotient with a metric on the same set,
see the author's
\cite[Proposition~2.12]{Ishiki2023FractalDimensions}.
We use the following correspondence form,
which also compares two metric quotients.

\begin{proposition}\label{lem:pseudometric-comparison}
\paraid{p-03dd7915a29a}
Let
$\BaseCarrier$
and
$\ComparisonCarrier$
be nonempty compact spaces.
Let
$\Correspondence\subset\BaseCarrier\times\ComparisonCarrier$
be a correspondence.
For continuous pseudometrics
$\Pseudometric$
on
$\BaseCarrier$
and
$\ComparisonPseudometric$
on
$\ComparisonCarrier$,
put
\begin{equation}\label{eq:correspondence-pseudometric-error}
 \PseudometricError_\Correspondence(\Pseudometric,\ComparisonPseudometric)
 =\sup_{(\BasePoint_0,\ComparisonPoint_0),(\BasePoint_1,\ComparisonPoint_1)\in\Correspondence}
 |\Pseudometric(\BasePoint_0,\BasePoint_1)
       -\ComparisonPseudometric(\ComparisonPoint_0,\ComparisonPoint_1)|.
\end{equation}
\begin{enumerate}[label=\textup{(Q\arabic*)},ref=\textup{(Q\arabic*)}]
\item\label{item:quotient-1}
\paraid{p-31226097263c}
The quotient spaces satisfy
\begin{equation}\label{eq:pseudometric-correspondence-bound}
 \GHDistance(\MetricQuotient{\BaseCarrier}{\Pseudometric},
              \MetricQuotient{\ComparisonCarrier}{\ComparisonPseudometric})
 \leq\tfrac12\PseudometricError_\Correspondence(\Pseudometric,\ComparisonPseudometric).
\end{equation}
\item\label{item:quotient-2}
\paraid{p-ef3a293af6ca}
If
$\Pseudometric_0,\Pseudometric_1$
are continuous pseudometrics on
$\BaseCarrier$
and
$\ComparisonPseudometric_0,\ComparisonPseudometric_1$
are continuous pseudometrics on
$\ComparisonCarrier$,
then
\begin{equation}\label{eq:pseudometric-uniform-comparison}
 \bigl|\norm{\Pseudometric_0-\Pseudometric_1}_\infty
       -\norm{\ComparisonPseudometric_0-\ComparisonPseudometric_1}_\infty\bigr|
 \leq\PseudometricError_\Correspondence(\Pseudometric_0,\ComparisonPseudometric_0)
       +\PseudometricError_\Correspondence(\Pseudometric_1,\ComparisonPseudometric_1).
\end{equation}
\item\label{item:quotient-3}
\paraid{p-b2b1a301c639}
Let the pseudometrics satisfy the hypotheses in \ref{item:quotient-2}
and let
$\HomotopyTime\in[0,1]$.
Then
\begin{equation}\label{eq:pseudometric-interpolation}
 \PseudometricError_\Correspondence((1-\HomotopyTime)\Pseudometric_0+\HomotopyTime\Pseudometric_1,
                   (1-\HomotopyTime)\ComparisonPseudometric_0+\HomotopyTime\ComparisonPseudometric_1)
 \leq(1-\HomotopyTime)\PseudometricError_\Correspondence(\Pseudometric_0,\ComparisonPseudometric_0)
       +\HomotopyTime\PseudometricError_\Correspondence(\Pseudometric_1,\ComparisonPseudometric_1).
\end{equation}
\item\label{item:quotient-4}
\paraid{p-45401a311836}
For continuous pseudometrics
$\Pseudometric_0,\Pseudometric_1$
on a nonempty compact space
$\BaseCarrier$,
\begin{equation}\label{eq:pseudometric-bound}
 \GHDistance(\MetricQuotient{\BaseCarrier}{\Pseudometric_0},\MetricQuotient{\BaseCarrier}{\Pseudometric_1})
 \leq\tfrac12\norm{\Pseudometric_0-\Pseudometric_1}_\infty.
\end{equation}
\end{enumerate}
\end{proposition}

\begin{proof}
\textbf{Proof of \ref{item:quotient-1}. Correspondence estimate.}
\paraid{p-04c1d56a107a}
Let
$\QuotientMap\colon\BaseCarrier\to\BaseCarrier/{\sim_\Pseudometric}$
and
$\SecondQuotientMap\colon\ComparisonCarrier\to\ComparisonCarrier/{\sim_\ComparisonPseudometric}$
be the quotient maps.
The set
\begin{equation}\label{eq:quotient-correspondence}
 \{(\QuotientMap(\BasePoint),\SecondQuotientMap(\ComparisonPoint))
       \mid(\BasePoint,\ComparisonPoint)\in\Correspondence\}
\end{equation}
is a correspondence because both quotient maps and both coordinate
projections of
$\Correspondence$
are surjective.
For two elements of the correspondence in
\eqref{eq:quotient-correspondence},
definition \eqref{eq:correspondence-pseudometric-error} implies
\begin{equation}\label{eq:quotient-correspondence-error}
 |\Pseudometric(\BasePoint_0,\BasePoint_1)
   -\ComparisonPseudometric(\ComparisonPoint_0,\ComparisonPoint_1)|
 \leq\PseudometricError_\Correspondence(\Pseudometric,\ComparisonPseudometric).
\end{equation}
Applying \Cref{thm:gh-correspondence} to \eqref{eq:quotient-correspondence-error}
proves \eqref{eq:pseudometric-correspondence-bound}.

\textbf{Proof of \ref{item:quotient-2}. Uniform comparison.}
\paraid{p-1e5ed86ba147}
Fix
$(\BasePoint_0,\ComparisonPoint_0),(\BasePoint_1,\ComparisonPoint_1)\in\Correspondence$.
The triangle inequality on the real line implies
\begin{equation}\label{eq:pseudometric-pointwise-comparison}
 |\Pseudometric_0(\BasePoint_0,\BasePoint_1)
    -\Pseudometric_1(\BasePoint_0,\BasePoint_1)|
 \leq|\ComparisonPseudometric_0(\ComparisonPoint_0,\ComparisonPoint_1)
         -\ComparisonPseudometric_1(\ComparisonPoint_0,\ComparisonPoint_1)|
 +\PseudometricError_\Correspondence(\Pseudometric_0,\ComparisonPseudometric_0)
       +\PseudometricError_\Correspondence(\Pseudometric_1,\ComparisonPseudometric_1).
\end{equation}
Every pair in
$\BaseCarrier\times\BaseCarrier$
has a corresponding pair in
$\ComparisonCarrier\times\ComparisonCarrier$.
Taking the supremum in \eqref{eq:pseudometric-pointwise-comparison},
we obtain
\begin{equation}\label{eq:pseudometric-sup-comparison}
 \norm{\Pseudometric_0-\Pseudometric_1}_\infty
 \leq\norm{\ComparisonPseudometric_0-\ComparisonPseudometric_1}_\infty
      +\PseudometricError_\Correspondence(\Pseudometric_0,\ComparisonPseudometric_0)
      +\PseudometricError_\Correspondence(\Pseudometric_1,\ComparisonPseudometric_1).
\end{equation}
Applying \eqref{eq:pseudometric-sup-comparison} after interchanging the carriers
proves \eqref{eq:pseudometric-uniform-comparison}.

\textbf{Proof of \ref{item:quotient-3}. Interpolation.}
\paraid{p-94f3dbc66479}
Fix
$(\BasePoint_0,\ComparisonPoint_0),(\BasePoint_1,\ComparisonPoint_1)\in\Correspondence$.
The difference of the interpolated pseudometrics at these pairs is
\begin{equation}\label{eq:pseudometric-pointwise-interpolation}
 (1-\HomotopyTime)
   \bigl(\Pseudometric_0(\BasePoint_0,\BasePoint_1)
       -\ComparisonPseudometric_0(\ComparisonPoint_0,\ComparisonPoint_1)\bigr)
 +\HomotopyTime
   \bigl(\Pseudometric_1(\BasePoint_0,\BasePoint_1)
       -\ComparisonPseudometric_1(\ComparisonPoint_0,\ComparisonPoint_1)\bigr).
\end{equation}
The absolute value of \eqref{eq:pseudometric-pointwise-interpolation} is bounded by the weighted sum of the two errors,
since both weights are nonnegative.
Taking the supremum and using \eqref{eq:correspondence-pseudometric-error}
proves \eqref{eq:pseudometric-interpolation}.

\textbf{Proof of \ref{item:quotient-4}. A common carrier.}
\paraid{p-95c5873e8f8e}
Use
$\Correspondence=\{(\BasePoint,\BasePoint)\mid\BasePoint\in\BaseCarrier\}$.
By \eqref{eq:correspondence-pseudometric-error},
\begin{equation}\label{eq:pseudometric-diagonal-error}
 \PseudometricError_\Correspondence(\Pseudometric_0,\Pseudometric_1)
 =\norm{\Pseudometric_0-\Pseudometric_1}_\infty.
\end{equation}
Substituting \eqref{eq:pseudometric-diagonal-error} into
\eqref{eq:pseudometric-correspondence-bound} proves \eqref{eq:pseudometric-bound}.
\end{proof}

% Retired paragraph IDs.
% \paraid{p-1248fc16354c}
% \paraid{p-83a90d959336}
